Considerando o espalhamento elástico elétron-múon, os tensores leptônico e hadrônico são dados por
    \begin{equation*}
        L^{\alpha\beta} = \dfrac{1}{2} \sum_{s,s'} 
        \bar{u}_{e}(k',s') \gamma^{\alpha} u_{e}(k,s)
        \bar{u}_{e}(k,s) \gamma^{\beta} u_{e}(k',s')
    \end{equation*}
e
    \begin{equation*}
        W_{\alpha\beta} = \dfrac{1}{2} \qty[\sum_{S,S'} 
        \bar{u}_{\mu}(p',S') \gamma_{\alpha} u_{\mu}(p,S)
        \bar{u}_{\mu}(p,S) \gamma_{\beta} u_{\mu}(p',S')]\delta^{4}(2p\cdot q + q^{2})
    \end{equation*}

Aqui consideramos $\bar{u}_{e}$ e $u_{e}$ como os espinores do elétron, $\bar{u}_{\mu}$ e $u_{\mu}$ os espinores do múon, $k$ e $k'$ os momentos iniciais e finais do elétron, $p$ e $p'$ os momentos iniciais e finais do múon, $q$ o momento transferido (como um fóton virtual), e $s,s',S,S'$ os índices de spin. O diagrama abaixo descreve o espalhamento em questão:
\begin{center}
    \begin{tikzpicture}
        \begin{feynman}
            \vertex (i1) {$e$};
            \vertex [right=3cm of i1, dot] (i2);
            \vertex [above right=1cm and 2.7cm of i2] (i3) {$e$};

            \vertex [below right=3cm and 1.5cm of i1] (i4) {$\mu$};
            \vertex [above right=1cm and 3cm of i4, dot] (i5);
            \vertex [right=3cm of i5] (i6) {$\mu$};

            \diagram*{
                (i1) -- [fermion, rmomentum'={[arrow shorten=0.7]$k,s$}] (i2) -- [fermion, rmomentum'={[arrow shorten=0.7]$k',s'$}] (i3),
                (i2) -- [boson, rmomentum'={[arrow shorten=0.7]$q$}] (i5),
                (i4) -- [fermion, rmomentum'={[arrow shorten=0.7]$p,S$}] (i5) -- [fermion, rmomentum'={[arrow shorten=0.7]$p',S'$}] (i6),
            };
        \end{feynman}
    \end{tikzpicture}
\end{center}

Fica claro então, por conservação de momento, que $q = k - k' = p' - p$. É visível que os tensores $L^{\alpha\beta}$ e $W_{\alpha\beta}$ possuem uma estrutura semelhante, então tratar somente uma delas é o suficiente para o cálculo da contração $L^{\alpha\beta}W_{\alpha\beta}$. A estrutura em questão é dada por
    \begin{equation*}
        \dfrac{1}{2} \sum_{s,s'} 
        \bar{u}(k',s') \gamma^{\alpha} u(k,s)
        \bar{u}(k,s) \gamma^{\beta} u(k',s')
    \end{equation*}
onde, para simplificar a notação, os índices $e$ e $\mu$ foram removidos. A expressão acima pode ser reescrita em componentes, isto é
    \begin{equation*}
        \dfrac{1}{2} \sum_{s,s'} \sum_{\tau,\sigma,\rho,\eta} 
        \bar{u}_{\tau}(k',s') (\gamma^{\alpha})_{\tau\sigma} u_{\sigma}(k,s)
        \bar{u}_{\rho}(k,s) (\gamma^{\beta})_{\rho\eta} u_{\eta}(k',s')
    \end{equation*}

No exercício 5.(b) da Lista 1, mostramos que
    \begin{equation*}
        \sum_{s} u_{\sigma}(k,s) \bar{u}_{\rho}(k,s) = (\slashed{k} + m)_{\sigma\rho}
    \end{equation*}

Portanto a soma em $s$ pode ser feita de modo a obtermos
    \begin{equation*}
        \dfrac{1}{2} \sum_{s'} \sum_{\tau,\sigma,\rho,\eta} 
        \bar{u}_{\tau}(k',s') (\gamma^{\alpha})_{\tau\sigma} 
        (\slashed{k} + m)_{\sigma\rho}
        (\gamma^{\beta})_{\rho\eta} u_{\eta}(k',s')
    \end{equation*}

Tratando-se de componentes matriciais, os termos podem ser rearranjados arbitrariamente, em particular a troca
    \begin{equation*}
        \dfrac{1}{2} \sum_{s'} \sum_{\tau,\sigma,\rho,\eta} 
        u_{\eta}(k',s') \bar{u}_{\tau}(k',s') 
        (\gamma^{\alpha})_{\tau\sigma} 
        (\slashed{k} + m)_{\sigma\rho}
        (\gamma^{\beta})_{\rho\eta}
    \end{equation*}

Performando então a soma em $s'$, obtemos
    \begin{equation*}
        \dfrac{1}{2} \sum_{\tau,\sigma,\rho,\eta} 
        (\slashed{k}' + m)_{\eta\tau} 
        (\gamma^{\alpha})_{\tau\sigma} 
        (\slashed{k} + m)_{\sigma\rho}
        (\gamma^{\beta})_{\rho\eta}
    \end{equation*}

Sabendo então que $\displaystyle\sum_{a,b}A_{ab}B_{ba} = \text{Tr}(AB)$, é imediato escrever
    \begin{equation*}
        \dfrac{1}{2} \sum_{\tau,\sigma,\rho,\eta} 
        (\slashed{k}' + m)_{\eta\tau} 
        (\gamma^{\alpha})_{\tau\sigma} 
        (\slashed{k} + m)_{\sigma\rho}
        (\gamma^{\beta})_{\rho\eta} = 
        \dfrac{1}{2}\text{Tr}\qty[
            (\slashed{k}' + m) 
            \gamma^{\alpha}
            (\slashed{k} + m)
            \gamma^{\beta}
        ]
    \end{equation*}

Sendo então $m$ a massa do elétron e $M$ a massa do múon, os tensores podem ser escritos sob a forma
    \begin{equation*}
        L^{\alpha\beta} = \dfrac{1}{2} \text{Tr}\qty[
            (\slashed{k}' + m)
            \gamma^{\alpha}
            (\slashed{k} + m)
            \gamma^{\beta}
        ]
    \end{equation*}
e
    \begin{equation*}
        W_{\alpha\beta} = \dfrac{1}{2} \text{Tr}\qty[
            (\slashed{p}' + M)
            \gamma_{\alpha}
            (\slashed{p} + M)
            \gamma_{\beta}
        ]\delta^{4}(2p\cdot q + q^{2})
    \end{equation*}

Tratando apenas o tensor leptônico, temos
    \begin{align*}
        L^{\alpha\beta} &\eq \dfrac{1}{2} \text{Tr}\qty[
            (\gamma^{\mu}k_{\mu}' + m)
            \gamma^{\alpha}
            (\gamma^{\nu}k_{\nu} + m)
            \gamma^{\beta}
        ] \\
        &\eq \dfrac{1}{2} \text{Tr}\qty[
            (\gamma^{\mu}\gamma^{\alpha}k_{\mu}' + \gamma^{\alpha}m)
            (\gamma^{\nu}\gamma^{\beta}k_{\nu} + \gamma^{\beta}m)
        ] \\
        &\eq \dfrac{1}{2} \text{Tr}\qty(
            \gamma^{\mu}\gamma^{\alpha}\gamma^{\nu}\gamma^{\beta} k_{\mu}'k_{\nu} +
            m\gamma^{\mu}\gamma^{\alpha}\gamma^{\beta} k_{\mu}' +
            m\gamma^{\alpha}\gamma^{\nu}\gamma^{\beta} k_{\nu} +
            m^{2}\gamma^{\alpha}\gamma^{\beta}
         ) \\
        &\eq \dfrac{1}{2} \qty[
            \text{Tr}(\gamma^{\mu}\gamma^{\alpha}\gamma^{\nu}\gamma^{\beta}) k_{\mu}'k_{\nu} + 
            m\text{Tr}(\gamma^{\mu}\gamma^{\alpha}\gamma^{\beta}) k_{\mu}' + 
            m\text{Tr}(\gamma^{\alpha}\gamma^{\nu}\gamma^{\beta}) k_{\nu} + 
            m^{2}\text{Tr}(\gamma^{\alpha}\gamma^{\beta}) 
        ]
    \end{align*}

Utilizando as propriedades do traço:
    \begin{align*}
        &\text{Tr}(\gamma^{\mu}\gamma^{\alpha}\gamma^{\nu}\gamma^{\beta}) = 4(g^{\mu\alpha}g^{\nu\beta} - g^{\mu\nu}g^{\alpha\beta} + g^{\mu\beta}g^{\alpha\nu}) \\
        &\text{Tr}(\gamma^{\mu}\gamma^{\alpha}\gamma^{\beta}) = \text{Tr}(\gamma^{\alpha}\gamma^{\nu}\gamma^{\beta}) = 0 \\
        &\text{Tr}(\gamma^{\alpha}\gamma^{\beta}) = 4g^{\alpha\beta}
    \end{align*}

Reduzimos a expressão a
    \begin{align*}
        L^{\alpha\beta} &\eq 2(g^{\mu\alpha}g^{\nu\beta} - g^{\mu\nu}g^{\alpha\beta} + g^{\mu\beta}g^{\alpha\nu}) k_{\mu}'k_{\nu} + 2 m^{2} g^{\alpha\beta} \\
        &\eq 2[k^{\prime\alpha}k^{\beta} - g^{\alpha\beta}(k'\cdot k) + k^{\prime\beta}k^{\alpha}] + 2 m^{2} g^{\alpha\beta} \\
        &\eq 2\qty[
            k^{\prime\alpha}k^{\beta} +
            k^{\prime\beta}k^{\alpha} - 
            g^{\alpha\beta}\qty(
                k'\cdot k - m^{2}
            )
        ]
    \end{align*}

Segue que o tensor hadrônico fica
    \begin{equation*}
        W_{\alpha\beta} = 2\Big[
            p_{\alpha}'p_{\beta} +
            p_{\beta}'p_{\alpha} - 
            g_{\alpha\beta}\qty(
                p'\cdot p - M^{2}
            )
        \Big]\delta^{4}(2p\cdot q + q^{2})
    \end{equation*}

Tendo então estas formas, a contração pode ser feita:
    \begin{align*}
        L^{\alpha\beta}W_{\alpha\beta} &\eq 4\qty[
            k^{\prime\alpha}k^{\beta} +
            k^{\prime\beta}k^{\alpha} - 
            g^{\alpha\beta}\qty(
                k'\cdot k - m^{2}
            )
        ]\Big[
            p_{\alpha}'p_{\beta} +
            p_{\beta}'p_{\alpha} - 
            g_{\alpha\beta}\qty(
                p'\cdot p - M^{2}
            )
        \Big] \times \\
        &\noeq \times \delta^{4}(2p\cdot q + q^{2}) \\
        &\eq 4 \Big[
            (k'\cdot p')(k\cdot p) +
            (k'\cdot p)(k\cdot p') -
            (k'\cdot k)\qty(
                p'\cdot p - M^{2}
            ) +
            (k'\cdot p)(k\cdot p') + \\
        &\noeq 
            (k'\cdot p')(k\cdot p) -
            (k'\cdot k)\qty(
                p'\cdot p - M^{2}
            ) -
            (p'\cdot p)\qty(
                k'\cdot k - m^{2}
            ) -
            (p'\cdot p)\qty(
                k'\cdot k - m^{2}
            ) + \\
        &\noeq 
            4\qty(
                k'\cdot k - m^{2}
            )\qty(
                p'\cdot p - M^{2}
            )
        \Big]\delta^{4}(2p\cdot q + q^{2}) \\
        &\eq 4\Big[
            2 (k'\cdot p')(k\cdot p) + 
            2 (k'\cdot p)(k\cdot p') -
            2 (k'\cdot k)\qty(
                p'\cdot p - M^{2}
            ) -
            2 (p'\cdot p)\qty(
                k'\cdot k - m^{2}
            ) + \\
        &\noeq +
            4 (k'\cdot k)\qty(
                p'\cdot p - M^{2}
            ) - 
            4 m^{2}\qty(
                p'\cdot p - M^{2}
            )
        \Big]\delta^{4}(2p\cdot q + q^{2}) \\
        &\eq 4\Big[
            2 (k'\cdot p')(k\cdot p) + 
            2 (k'\cdot p)(k\cdot p') +
            2 (k'\cdot k)\qty(
                p'\cdot p - M^{2}
            ) -
            2 (p'\cdot p)\qty(
                k'\cdot k - m^{2}
            ) + \\
        &\noeq -
            4 m^{2}\qty(
                p'\cdot p - M^{2}
            )
        \Big]\delta^{4}(2p\cdot q + q^{2}) \\
        &\eq 4\Big[
            2(k'\cdot p')(k\cdot p) +
            2(k'\cdot p)(k\cdot p') +
            2(k'\cdot k)(p'\cdot p) - 
            2M^{2}(k'\cdot k) - 
            2(k'\cdot k)(p'\cdot p) + \\
        &\noeq +
            2m^{2}(p'\cdot p) - 
            4m^{2}(p'\cdot p) + 
            4m^{2}M^{2}
        \Big]\delta(2p\cdot q + q^{2}) \\
        &\eq 8\qty[
            (k'\cdot p')(k\cdot p) +
            (k'\cdot p)(k\cdot p') -
            M^{2}(k'\cdot k) -
            m^{2}(p'\cdot p) +
            2m^{2}M^{2}
        ]\delta(2p\cdot q + q^{2}) 
    \end{align*}

Partindo pra física do problema, ao considerarmos o referencial do laboratório (múon em repouso no estado inicial), podemos escrever
    \begin{equation*}
        k^{\mu} = (E,\vb{k}) \qquad 
        k^{\prime\mu} = (E',\vb{k}') \qquad 
        p^{\mu} = (M,\vb{0}) \qquad 
        p^{\prime\mu} = (\mathcal{E},\vb{p}')
    \end{equation*}

Como estamos em altas energias ($E,E'\gg m$), a massa do elétron pode ser desconsiderada e pela relação de energia-momento, valem as igualdades $\abs{\vb{k}} = E$ e $\abs{\vb{k}'} = E'$. Então desconsiderando qualquer termo que envolva a massa do elétron, a contração tensorial fica
    \begin{equation*}
        L^{\alpha\beta}W_{\alpha\beta} = 8\qty[
            (k'\cdot p')(k\cdot p) +
            (k'\cdot p)(k\cdot p') -
            M^{2}(k'\cdot k)
        ]\delta(2p\cdot q + q^{2}) 
    \end{equation*}

Na conservação de momento, $p' = p + q$ e $q = k - k'$, então
    \begin{align*}
        k'\cdot p' = k'\cdot (p + q) &= k'\cdot p + k'\cdot q &
        k\cdot p' = k\cdot (p + q) &= k\cdot p + k\cdot q \\
        &= k'\cdot p + k'\cdot(k - k') & &= k\cdot p + k\cdot (k - k') \\
        &= k'\cdot p + k'\cdot k - (k')^{2} & &= k\cdot p + k^{2} - k\cdot k'
    \end{align*}

Por estamos desconsiderando a massa do elétron, $k^{2} = (k')^{2} = m^{2} \approx 0$, então temos
    \begin{align*}
        k'\cdot p' &= k'\cdot p + k'\cdot k & k\cdot p' = k\cdot p - k'\cdot k
    \end{align*}

Além disso, podemos ver que
    \begin{equation*}
        q^{2} = (k - k')^{2} = k^{2} + (k')^{2} - 2(k'\cdot k) \approx -2(k'\cdot k)
    \end{equation*}

Segue que
    \begin{align*}
        L^{\alpha\beta}W_{\alpha\beta} &\eq 8\qty[
            \qty(k'\cdot p - \dfrac{q^{2}}{2}) (k\cdot p) +
            (k'\cdot p)\qty(k\cdot p + \dfrac{q^{2}}{2}) + 
            \dfrac{M^{2}q^{2}}{2}
        ]\delta^{4}(2p\cdot q + q^{2}) \\
        &\eq 8\qty[
            (k'\cdot p)(k\cdot p) -
            \dfrac{q^{2}}{2}(k\cdot p) +
            (k'\cdot p)(k\cdot p) +
            \dfrac{q^{2}}{2}(k'\cdot p) +
            \dfrac{M^{2}q^{2}}{2}
        ]\delta^{4}(2p\cdot q + q^{2}) \\
        &\eq 8\qty[
            2(k'\cdot p)(k\cdot p) - 
            \dfrac{q^{2}}{2}(k\cdot p - k'\cdot p) +
            \dfrac{M^{2}q^{2}}{2}
        ]\delta^{4}(2p\cdot q + q^{2})
    \end{align*}


No referencial do laboratório, temos
    \begin{align*}
        (k'\cdot p) &= (E',\vb{k}')\cdot(M,\vb{0}) = ME' &
        (k\cdot p) &= (E,\vb{k})\cdot(M,\vb{0}) = ME
    \end{align*}

Então
    \begin{align*}
        L^{\alpha\beta}W_{\alpha\beta} &\eq 8\qty[
            2M^{2}E'E - 
            \dfrac{Mq^{2}}{2}(E - E') +
            \dfrac{M^{2}q^{2}}{2}
        ]\delta^{4}(2p\cdot q + q^{2})
    \end{align*}

Note que $q^{2}$ pode ser reescrito em termos de um ângulo $\theta$ formado pelos vetores $\vb{k}$ e $\vb{k}'$, isto é
    \begin{align*}
        q^{2} &\eq -2(k'\cdot k) \\
        &\eq -2(E'E - \vb{k}'\cdot\vb{k}) \\
        &\eq -2\qty(E'E - \abs{\vb{k}'}\abs{\vb{k}}\cos\theta) \\
        &\eq -2\qty(E'E - E'E\cos\theta) \\
        &\eq -2E'E(1 - \cos\theta) \\
        &\eq -4E'E\sin^{2}\qty(\dfrac{\theta}{2})
    \end{align*}

Portanto
    \begin{align*}
        L^{\alpha\beta}W_{\alpha\beta} &\eq 8\qty[
            2M^{2}E'E + 
            2ME'E(E-E')\sin^{2}\qty(\dfrac{\theta}{2}) -
            2M^{2}E'E\sin^{2}\qty(\dfrac{\theta}{2})
        ]\delta^{2}(2p\cdot q + q^{2}) \\
        &\eq 16 M^{2}E'E\qty[
            1 -
            \sin^{2}\qty(\dfrac{\theta}{2}) +
            \dfrac{(E - E')}{M}\sin^{2}\qty(\dfrac{\theta}{2})
        ]\delta^{2}(2p\cdot q + q^{2}) \\
        &\eq 16 M^{2}E'E\qty[
            \cos^{2}\qty(\dfrac{\theta}{2}) +
            \dfrac{(E - E')}{M}\sin^{2}\qty(\dfrac{\theta}{2})
        ]\delta^{4}(2p\cdot q + q^{2})
    \end{align*}

Podemos escrever
    \begin{align*}
        E - E' = \dfrac{1}{M}(ME - ME') = \dfrac{1}{M}(p\cdot k - p\cdot k') = \dfrac{p\cdot(k - k')}{M} = \dfrac{p\cdot q}{M}
    \end{align*}

Então definindo $\nu \coloneqq \dfrac{p\cdot q}{M} = E - E'$ e $q^{2} \equiv -Q^{2}$ (momento transferido), concluímos que a contração entre os tensores, no referencial do laboratório, assume a forma
\begin{answer}\label{eq: answer 6}
    L^{\alpha\beta}W_{\alpha\beta} = 16M^{2}E'E\qty[
        \cos^{2}\qty(\dfrac{\theta}{2}) + 
        \dfrac{\nu}{M}\sin^{2}\qty(\dfrac{\theta}{2})
    ]\delta^{4}(2M\nu - Q^{2})
\end{answer}

\endex